% ============================================================================
% ESL_Paper_v12_part2.tex
% ============================================================================
% This file contains the remaining sections of the ESL paper.
% It is included via % ============================================================================
% ESL_Paper_v12_part2.tex
% ============================================================================
% This file contains the remaining sections of the ESL paper.
% It is included via % ============================================================================
% ESL_Paper_v12_part2.tex
% ============================================================================
% This file contains the remaining sections of the ESL paper.
% It is included via % ============================================================================
% ESL_Paper_v12_part2.tex
% ============================================================================
% This file contains the remaining sections of the ESL paper.
% It is included via \input{ESL_Paper_v12_part2} in the main document.
% ============================================================================

% TODO: Add remaining sections here
% Suggested sections:
% - Section 8: Implementation Guidelines
% - Section 9: Discussion and Limitations
% - Section 10: Future Work
% - Section 11: Conclusion
% - References (using \bibliography{references})

\section{Conclusion}

The Empirical Stability of Language (ESL) framework provides a systematic approach to ensuring epistemic integrity in assertive communication. By operationalizing the principle that claim strength must not exceed evidence strength ($K \leq M$), ESL enables calibrated scientific communication.

Our key finding---that LLMs already possess the knowledge required for epistemic calibration---suggests that expensive retraining may be unnecessary for achieving calibrated output. Instead, systematic activation through structured prompts can unlock latent competence.

ESL thus represents a paradigm shift: from knowledge injection to knowledge extraction, from prompting as art to system configuration as engineering.

% Bibliography placeholder
% \bibliographystyle{apalike}
% \bibliography{references}

\end{document}
 in the main document.
% ============================================================================

% TODO: Add remaining sections here
% Suggested sections:
% - Section 8: Implementation Guidelines
% - Section 9: Discussion and Limitations
% - Section 10: Future Work
% - Section 11: Conclusion
% - References (using \bibliography{references})

\section{Conclusion}

The Empirical Stability of Language (ESL) framework provides a systematic approach to ensuring epistemic integrity in assertive communication. By operationalizing the principle that claim strength must not exceed evidence strength ($K \leq M$), ESL enables calibrated scientific communication.

Our key finding---that LLMs already possess the knowledge required for epistemic calibration---suggests that expensive retraining may be unnecessary for achieving calibrated output. Instead, systematic activation through structured prompts can unlock latent competence.

ESL thus represents a paradigm shift: from knowledge injection to knowledge extraction, from prompting as art to system configuration as engineering.

% Bibliography placeholder
% \bibliographystyle{apalike}
% \bibliography{references}

\end{document}
 in the main document.
% ============================================================================

% TODO: Add remaining sections here
% Suggested sections:
% - Section 8: Implementation Guidelines
% - Section 9: Discussion and Limitations
% - Section 10: Future Work
% - Section 11: Conclusion
% - References (using \bibliography{references})

\section{Conclusion}

The Empirical Stability of Language (ESL) framework provides a systematic approach to ensuring epistemic integrity in assertive communication. By operationalizing the principle that claim strength must not exceed evidence strength ($K \leq M$), ESL enables calibrated scientific communication.

Our key finding---that LLMs already possess the knowledge required for epistemic calibration---suggests that expensive retraining may be unnecessary for achieving calibrated output. Instead, systematic activation through structured prompts can unlock latent competence.

ESL thus represents a paradigm shift: from knowledge injection to knowledge extraction, from prompting as art to system configuration as engineering.

% Bibliography placeholder
% \bibliographystyle{apalike}
% \bibliography{references}

\end{document}
 in the main document.
% ============================================================================

% TODO: Add remaining sections here
% Suggested sections:
% - Section 8: Implementation Guidelines
% - Section 9: Discussion and Limitations
% - Section 10: Future Work
% - Section 11: Conclusion
% - References (using \bibliography{references})

\section{Conclusion}

The Empirical Stability of Language (ESL) framework provides a systematic approach to ensuring epistemic integrity in assertive communication. By operationalizing the principle that claim strength must not exceed evidence strength ($K \leq M$), ESL enables calibrated scientific communication.

Our key finding---that LLMs already possess the knowledge required for epistemic calibration---suggests that expensive retraining may be unnecessary for achieving calibrated output. Instead, systematic activation through structured prompts can unlock latent competence.

ESL thus represents a paradigm shift: from knowledge injection to knowledge extraction, from prompting as art to system configuration as engineering.

% Bibliography placeholder
% \bibliographystyle{apalike}
% \bibliography{references}

\end{document}
